\Needspace{6\baselineskip}\Article{331}
In Section IV we explained in detail how to distinguish the primes of which a given number \(r\) (which we may assume is divisible by no square) is a residue from those of which it is a nonresidue, or the divisors of the expression \(xx-r\) from its non-divisors: all the former are contained in certain formulas of the kind \(rz+a,\allowbreak \ rz+b\), etc., or of the kind \(4rz+a,\allowbreak \ 4rz+b\), etc., and the latter in other similar formulas. Whenever \(r\) is a very small number, these formulas allow the exclusions to be made
\SourceLeaf{409}
very conveniently; for example, we must exclude all numbers of the form \(4z+3\) when \(r=-1\); all numbers of the forms \(8z+3\) and \(8z+5\) when \(z=2\), etc. But we cannot always find such residues of the proposed number \(M\), and applying the formulas for a large value of \(r\) is not convenient enough. It is therefore an enormous advantage, and remarkably eases the work of exclusion, to have a \textit{table} already constructed for a sufficiently large collection of numbers \((r)\), both positive and negative, divisible by no square, distinguishing the primes of which each of these \((r)\) is a residue from those of which it is a nonresidue. Such a table can be arranged like the sample appended to this work and already described above; but to be sufficiently useful for the present purpose, the primes placed in the margin (the moduli) must be continued much further, say at least to \(1000\) or \(10000\). It is also much more convenient to include composite and negative numbers among the headings, although, as Section IV shows, this is not absolutely necessary. Such a table will be most convenient to use if its individual vertical columns are cut out and glued onto strips or rods (like Napier's), so that those needed in any given case, namely those corresponding to the numbers \(r,\allowbreak r',\allowbreak r''\), etc., residues of the proposed number to be factored, can be examined separately. When they are placed \textit{properly} beside the first column of the table (which gives the moduli), that is, so that the positions on the individual rods corresponding to the same prime in the first column are aligned with it, or lie on the same horizontal line, we can clearly recognize at a glance the primes that, after the exclusions with the residues \(r,\allowbreak r',\allowbreak r''\), remain in \(\Omega\): they are precisely those numbers in the first column that have a stroke beside them on \textit{all} the adjacent rods, and we must reject all those with a blank space beside them on \textit{any} rod. An example will make this sufficiently clear. If it is known from some source that the numbers \(-6,\allowbreak +13,\allowbreak -14,\allowbreak +17,\allowbreak +37,\allowbreak -53\) are residues of \(997331\), we must put together the first column (which in this case must extend to \(997\), the prime immediately below \(\surd{997331}\)) and the strips headed by the numbers \(-6,\allowbreak +13\), etc. Here is part of the resulting arrangement:
\SourceLeaf{410}
\[\renewcommand{\arraystretch}{1.05}\begin{array}{|r|c|c|c|c|c|c|}&-6&+13&-14&+17&+37&-53\\\hline 3&-&-&-&&-&-\\5&-&&-&&&\\7&-&&-&&-&\\11&-&&&&-&\\13&&-&-&-&&-\\17&&-&&-&&-\\19&&&-&-&&-\\23&&-&-&&&-\\&&&\text{etc.}&&&\\113&&-&-&&&-\\127&-&-&-&-&-&-\\131&-&-&-&&&\\&&&\text{etc.}&&&\end{array}\]

Just as a glance shows here that, among the primes \textit{included in this part of the arrangement}, only \(127\) remains after the exclusions with the residues \(-6,\allowbreak 13\), etc., in \(\Omega\), so the complete arrangement extended to \(997\) shows that \textit{absolutely} no other prime remains in \(\Omega\); on trying the division, we find that \(997331\) is indeed divisible by \(127\). In this way, therefore, that number has been resolved into the prime factors \(127\times7853\)\HistoricalNote{*)}{The author would gladly publish the fairly extensive apparatus of the table described here, which he had constructed for his own use, if the small number of people to whom it can be useful were sufficient to cover the costs of such an undertaking. If, meanwhile, any lover of arithmetic, having thoroughly understood the principles, wishes to construct such a table for himself by his own efforts, the author will take great pleasure in sharing with him by letter all the aids and devices.}.

This account also makes it sufficiently clear that residues which are not too large, or can at least be resolved into prime factors which are not too large, are particularly useful: the direct use of the auxiliary table does not extend beyond the numbers among its headings, and its indirect use covers only those that can be resolved into factors contained in the table.

\Needspace{6\baselineskip}\Article{332}
We will give three different methods of finding residues of a given number \(M\). Before explaining them, we make two observations that allow simpler residues to be derived from less suitable ones. \textit{First}, if a number \(akk\) divisible by the square \(kk\) (which is assumed coprime to \(M\)) is a residue of \(M\), then \(a\) is also a residue; for this reason, residues divisible by large
\SourceLeaf{411}
squares are just as useful as small ones, and we will assume that the square factors of all residues supplied by the following methods have been removed at once. \textit{Second}, if two or more numbers are residues, their product is also a residue. Combining this observation with the preceding one very often allows us to derive a very simple residue from several that are not all simple enough, provided they contain many common factors. For this reason, residues composed of many factors which are not too large are also very suitable, and it will be useful to resolve them all into their factors at once. The force of these observations will be better understood through examples and frequent use than through rules.

I. The simplest method, and the most convenient for those who have already gained some skill through frequent practice, consists in decomposing \(M\), or more generally any multiple of \(M\), in any way into two parts \(kM=a+b\) (whether both are positive or one positive and the other negative). Their product with its sign changed will be a residue of \(M\); for \(-ab\equiv aa\equiv bb\ (\mathrm{mod.}\ M)\), and hence \(-ab\,R\,M\). The numbers \(a,\allowbreak b\) should be chosen so that their product is divisible by a large square and the quotient is either small or at least resolvable into factors which are not too large; this can always be done without difficulty. We particularly recommend taking for \(a\) a square, or twice a square, or three times a square, etc., differing from the number \(M\) by a number which is either small or resolvable into convenient factors. For example, we find \[\begin{aligned}
997331&=999^2-2.5.67\\
&=994^2+5.11.13^2=2.706^2+3.17.3^2\\
&=3.575^2+11.31.4^2=3.577^2-7.13.4^2\\
&=3.578^2-7.19.37=11.299^2+2.3.5.29.4^2\\
&=11.301^2+5.12^2
\end{aligned}\] etc. This gives the following residues: \(2.5.67,\allowbreak \ -5.11,\allowbreak \ -2.3.17,\allowbreak \ -3.11.31,\allowbreak \ 3.7.13,\allowbreak \ 3.7.19.37,\allowbreak -2.3.5.11.29\); the last decomposition supplies the residue \(-5.11\), which we already have. Instead of the residues \(-3.11.31,\allowbreak \ -2.3.5.11.29\), we can take \(3.5.31,\allowbreak \ 2.3.29\), obtained by combining them with \(-5.11\).

\begin{samepage}
II. The second and third methods follow from the fact that, if two binary forms \((A,\allowbreak B,\allowbreak C),\allowbreak \ (A',\allowbreak B',\allowbreak C')\) of the same determinant \(M\), or \(-M\), or more generally \(\pm kM\), belong to the same genus, the numbers \(AA',\allowbreak \ AC',\allowbreak \ A'C\) are residues of \(kM\). This is readily seen from the fact that any characteristic number of one form, say \(m\), is also a characteristic number of the other, and hence \(mA,\allowbreak \ mC,\allowbreak \)
\SourceLeaf{412}
\(mA',\allowbreak \ mC'\) are all residues of \(kM\). Thus, if \((a,\allowbreak b,\allowbreak a')\) is a reduced form of positive determinant \(M\), or more generally \(kM\), and \((a',\allowbreak b',\allowbreak a''),\allowbreak \ (a'',\allowbreak b'',\allowbreak a''')\), etc., are forms in its period, and hence equivalent to it and all the more certainly contained in the same genus, the numbers \(aa',\allowbreak \ aa'',\allowbreak \ aa'''\), etc., will all be residues of \(M\). A large number of forms in such a period can be calculated very easily with the algorithm of Article 187; the simplest residues usually result from setting \(a=1\); those containing factors which are too large should be rejected. Here are the beginnings of the periods of the forms \((1,\allowbreak 998,\allowbreak -1327)\) and \((1,\allowbreak 1412,\allowbreak -918)\), whose determinants are \(997331,\allowbreak \ 1994662\): \[\begin{array}{rrr|rrr}\bigl(1,&998,&-1327\bigr)&\bigl(1,&1412,&-918\bigr)\\ \bigl(-1327,&329,&670\bigr)&\bigl(-918,&1342,&211\bigr)\\ \bigl(670,&341,&-1315\bigr)&\bigl(211,&1401,&-151\bigr)\\ \bigl(-1315,&974,&37\bigr)&\bigl(-151,&1317,&1723\bigr)\\ \bigl(37,&987,&-626\bigr)&\bigl(1723,&406,&-1062\bigr)\\ \bigl(-626,&891,&325\bigr)&\bigl(-1062,&656,&1473\bigr)\\ \bigl(325,&734,&-1411\bigr)&\bigl(1473,&817,&-901\bigr)\\ \bigl(-1411,&677,&382\bigr)&\bigl(-901,&985,&1137\bigr)\\ \bigl(382,&851,&-715\bigr)&\multicolumn{3}{c}{\mathrm{etc.}}\end{array}\] Thus the number \(997331\) has all the numbers \(-1327,\allowbreak 670\), etc., as residues; omitting those containing factors which are too large, we have the following: \(2.5.67,\allowbreak \ 37,\allowbreak \ 13,\allowbreak -17.83,\allowbreak \ -5.11.13,\allowbreak \ -2.3.17,\allowbreak \ -2.59,\allowbreak \ -17.53\). The residue \(2.5.67\), as well as \(-5.11\), which results from combining the third with the fifth, had already been found above.
\par\end{samepage}

III. If \(C\) is any class of forms of negative determinant \(-M\), or more generally \(-kM\), different from the principal class, and its period is \(2C,\allowbreak \ 3C\), etc. (Article 307), the classes \(2C,\allowbreak \ 4C\), etc., will belong to the principal genus, and \(3C,\allowbreak \ 5C\), etc., to the same genus as \(C\). Thus, if \((a,\allowbreak b,\allowbreak c)\) is a (simplest) form in \(C\), and \((a',\allowbreak b',\allowbreak c')\) a form in some class of that period, say \(nC\), either \(a'\) or \(aa'\) will be a residue of \(M\), according as \(n\) is even or odd (in the former case clearly also \(c'\), and in the latter \(ac',\allowbreak \ ca'\) and \(cc'\)). The period, that is, the simplest forms in its classes, can be worked out with remarkable ease when \(a\) is very small, especially in the case \(=3\), which can always be achieved when \(kM\equiv2\ (\mathrm{mod.}\ 3)\). Here is the beginning of the period of the class containing the form \((3,\allowbreak 1,\allowbreak 332444)\):
\SourceLeaf{413}
\[\begin{array}{crrr|crrr}C&\bigl(3,&1,&332444\bigr)&6C&\bigl(729.&-209,&1428\bigr)\\ 2C&\bigl(9,&-2,&110815\bigr)&7C&\bigl(476,&209,&2187\bigr)\\ 3C&\bigl(27,&7,&36940\bigr)&8C&\bigl(1027,&342,&1085\bigr)\\ 4C&\bigl(81,&34,&12327\bigr)&9C&\bigl(932,&-437,&1275\bigr)\\ 5C&\bigl(243,&34,&4109\bigr)&10C&\bigl(425,&12,&2347\bigr)\end{array}\] This gives the residues (rejecting the useless ones) \(3.476,\allowbreak \ 1027,\allowbreak \ 1085,\allowbreak \ 425\), or (removing square factors) \(3.7.17,\allowbreak \ 13.79,\allowbreak \ 5.7.31,\allowbreak \ 17\). By combining these suitably with the eight residues found in II, we easily obtain the following twelve: \(-2.3,\allowbreak \ 13,\allowbreak \ -2.7,\allowbreak 17,\allowbreak \ 37,\allowbreak \ -53,\allowbreak \ -5.11,\allowbreak \ 79,\allowbreak \ -83,\allowbreak \ -2.59,\allowbreak \ -2.5.31,\allowbreak \ 2.5.67\); the first six are the same as those we used in Article 331. We could have added the residues \(19\) and \(-29\) if we had also wished to use those found in I; the others obtained there already depend on those developed here.

\Needspace{6\baselineskip}\Article{333}
The \textsc{second method} of factoring a given number \(M\) comes from considering the values of an expression such as \(\surd{-D}\ (\mathrm{mod.}\ M)\), and rests on the following observations.

I. When \(M\) is a prime or a power of a prime (the prime being odd and not dividing \(D\)), \(-D\) will be a residue or a nonresidue of \(M\) according as \(M\) is contained in a form of divisors or in a form of non-divisors of \(xx+D\); in the former case the expression \(\surd{-D}\ (\mathrm{mod.}\ M)\) will have only two distinct values, which will be opposite.

II. But when \(M\) is composite, say \(=pp'p''\), etc., where \(p,\allowbreak p',\allowbreak p''\), etc., denote primes (distinct, odd, and not dividing \(D\)) or powers of such primes, \(-D\) will be a residue of \(M\) only when it is a residue of each of \(p,\allowbreak p',\allowbreak p''\), etc., that is, when all these numbers are contained in forms of divisors of \(xx+D\). If we denote the values of the expression \(\surd{-D}\) modulo \(p,\allowbreak p',\allowbreak p''\), etc., respectively, by \(\pm r,\allowbreak \ \pm r',\allowbreak \ \pm r''\), etc., all the values of the same expression modulo \(M\) arise by finding the numbers which, modulo \(p\), are \(\equiv r\) or \(\equiv-r\), modulo \(p'\), either \(\equiv r'\) or \(\equiv-r'\), etc. Their number will therefore be \(=2^\mu\), where \(\mu\) counts the numbers \(p,\allowbreak p',\allowbreak p''\), etc. Thus, if these values are \(R,\allowbreak -R,\allowbreak R',\allowbreak -R',\allowbreak R''\), etc., we automatically have \(R\equiv R\) modulo all of \(p,\allowbreak p',\allowbreak p''\), etc., but modulo none of them \(R\equiv-R\). Hence the greatest common
\SourceLeaf{414}
divisor of \(M\) and \(R-R\) will be \(M\), and \(1\) will be the greatest common divisor of \(M\) and \(R+R\). But two values which are neither identical nor opposite, such as \(R\) and \(R'\), must be congruent modulo one or more of the numbers \(p,\allowbreak p',\allowbreak p''\), etc., but not modulo all of them, and modulo the rest \(R\equiv-R'\); hence the product of the former will be the greatest common divisor of \(M\) and \(R-R'\), and the product of the latter the greatest common divisor of \(M\) and \(R+R'\). It follows easily that, if we calculate all the greatest common divisors of \(M\) and the differences between the individual values of the expression \(\surd{-D}\ (\mathrm{mod.}\ M)\) and any given value, their collection contains the numbers \(1,\allowbreak p,\allowbreak p',\allowbreak p''\), etc., and all products of two, three, etc., of these numbers. \textit{In this way, therefore, from} \textit{the values of that expression we can obtain the numbers \(p,\allowbreak p',\allowbreak p''\), etc.}

Moreover, the method of Article 327 reduces each of these values to the values of expressions of the form \(\frac mn\ (\mathrm{mod.}\ M)\), where the denominator \(n\) is coprime to \(M\); for the present purpose it is not even necessary to calculate these expressions themselves. For the greatest common divisor of \(M\) and the difference between \(R\) and \(R'\), which agree with \(\frac mn,\allowbreak \frac{m'}{n'}\), is clearly also the greatest common divisor of \(M\) and \(nn'(R-R')\). Equivalently, it is the greatest common divisor of \(M\) and \(mn'-m'n\), to which \(nn'(R-R')\) is congruent modulo \(M\).

\Needspace{12\baselineskip}\Article{334}
\enlargethispage{\baselineskip}
The preceding observations can be applied to our problem in two ways. The first not only decides whether the proposed number \(M\) is prime or composite, but in the latter case also supplies the factors themselves; the second has the advantage of usually allowing a quicker calculation, but sometimes fails to produce the factors themselves of composite numbers, which it likewise distinguishes completely from primes, unless it is repeated several times.

I. Find a negative number \(-D\) which is a quadratic residue of \(M\); for this purpose the methods given in Article 332 under I and II can be used. In itself, the choice of residue is arbitrary, and here, unlike the preceding method, we do not need \(D\) to be a small number; but the calculation will be shorter, the smaller the number of classes of binary forms in each properly primitive genus of determinant \(-D\). Hence residues among the \(65\) numbers of Article 303, if any present themselves, will be particularly suitable. Thus for \(M=997331\), of all the negative residues obtained above, \(-102\) would be the most
\SourceLeaf{415}
suitable. Find all the distinct values of the expression \(\surd{-D}\ (\mathrm{mod.}\ M)\); if only two result (opposite ones), \(M\) will certainly be either a prime or a power of a prime; if there are more, say \(2^\mu\), \(M\) will be composed of \(\mu\) distinct primes, or powers of primes, and these factors can be obtained by the method of the preceding article. Whether these factors are primes or powers of primes will be very easy to recognize in itself; furthermore, the very procedure by which the values of the expression \(\surd{-D}\) are found automatically indicates all primes some power of which divides \(M\). Namely, if \(M\) is divisible by the square of a prime \(\pi\), that calculation will certainly also have produced one or more representations of \(M\), \(M=amm+2bmn+cnn\), in which the greatest common divisor of the numbers \(m,\allowbreak n\) is \(\pi\) (indeed, because in this case \(-D\) is also a residue of \(\frac{M}{\pi\pi}\)). But when no representation has appeared in which \(m\) and \(n\) have a common divisor, this is a sure indication that \(M\) is divisible by no square, and hence all of \(p,\allowbreak p',\allowbreak p''\), etc., are primes.

Example. By the method given above we find four values of the expression \(\surd{-408}\ (\mathrm{mod.}\ 997331)\) agreeing with the values of \(\pm\frac{1664}{113},\allowbreak \ \pm\frac{2824}{3}\); the greatest common divisors of \(997331\) with \(3.1664-113.2824\) and \(3.1664+113.2824\), or with \(314120\) and \(324104\), will be \(7853\) and \(127\), whence \(997331=127.7853\), as above.

II. Take any negative number \(-D\) such that \(M\) is contained in a form of divisors of \(xx+D\). In itself, the choice of such a number is arbitrary, but for convenience we should particularly seek to make the number of classes in the genera of determinant \(-D\) as small as possible. Finding such a number presents no difficulty if we proceed by trial; for usually, among a considerable collection of numbers tried, \(M\) is contained in a form of divisors for about as many as in a form of non-divisors. It will therefore be most useful to begin the trial with the \(65\) numbers of Article 303 (starting with the largest), and if it should happen that none is suitable (which, generally speaking, occurs only once in \(16384\) cases), to proceed to others for which each genus contains two classes. — Then find the values of the expression \(\surd{-D}\ (\mathrm{mod.}\ M)\), and if any are found, deduce the factors of \(M\) from them in exactly the same way as above; but if no values appear, so that \(-D\) is a nonresidue of \(M\), then \(M\) certainly cannot be either a prime or a power of a prime. If, in this
\SourceLeaf{416}
case, the factors themselves are wanted, we must either repeat the same operation, choosing other values for \(D\), or turn to another method.

Thus, for example, on making the trial, we find that \(997331\) is contained in a form of non-divisors of \(xx+1848,\allowbreak \ xx+1365,\allowbreak \ xx+1320\), but in a form of divisors of \(xx+840\); the values of the expression \(\surd{-840}\ (\mathrm{mod.}\ 997331)\) give the expressions \(\pm\frac{1272}{163},\allowbreak \pm\frac{3288}{125}\), from which the same factors are deduced as before. — Anyone wanting further examples may consult Article 328: the first shows that \(5428681=307.17683\); the second, that \(4272943\) is a prime; the third, that \(997331\) is certainly composed of several primes.

\[\begin{array}{ccc}*& &*\\ &*&\end{array}\]

The limits of this work have allowed us to pursue here only the main points of the two methods of finding factors; we reserve a fuller investigation, together with further auxiliary tables and other aids, for another occasion.
\BookRule{0.22\textwidth}
\clearpage
